% --- Archon physics formalization source begin ---
% archon:physics
% archon:covers PhyXMiniProblems/problem_phyx_mini_0053.lean
% archon:source-report reports/phyx_mini/problem_phyx_mini_0053.source.json

\chapter{Physics problem phyx\_mini\_0053}
\label{ch:PhyXMiniProblems_problem_phyx_mini_0053}

\paragraph{Problem source.}
\#\# Physical scenario

A diverging lens with a focal length of 50 cm is placed 100 cm from a flower.

\#\# Question

What is the magnification?

\#\# Answer choices

- A: A: 0.49

- B: B: 0.33

- C: C: 0.14

- D: D: 0.84

\#\# Figure caption (auxiliary; use the image as primary evidence)

The image is a diagram illustrating the behavior of light passing through a diverging lens. It includes the following elements:

1. **Object**: Represented by an upright arrow on the left labeled "Object," positioned above the principal axis.

2. **Lens**: A biconcave (diverging) lens is drawn in the center, marked by the line labeled "a."

3. **Principal Axis**: A horizontal line running through the center of the lens.

4. **Focal Points**: There are focal points on either side of the lens. The left focal point is labeled with an "f," and another "f" is shown on the right, each at equal distances from the lens.

5. **Image**: Formed as an inverted arrow on the right side of the lens, below the principal axis, labeled "Image." It is located between the lens and the right focal point.

6. **Light Rays**: The diagram includes several light rays:
   - One ray passing through the center of the lens and continuing in a straight path.
   - Two other rays diverging after passing through the lens, appearing to originate from the focal point on the left side.

7. **Labels and Measurements**:
   - "s" and "s'" represent object and image distances from the lens, respectively.
   - The object height is labeled as "4" and the image height as "5."
   - A distance "b" is marked to the right of the lens, measuring 50 cm from the lens to a point on the axis.
   - A distance from the object to the lens is labeled "100."

8. **Angles and Directions**:
   - Arrows indicating the path and direction of the light rays and showing the relationships between object, lens, and image.

The diagram is likely for educational purposes, showing how a diverging lens affects the path of light and the nature of the image formed.

\#\# Dataset metadata

Geometrical Optics; Physical Model Grounding Reasoning, Spatial Relation Reasoning

\paragraph{Recorded answer/context.}
B: B: 0.33

\paragraph{Figure/image path.}
/root/proposal\_for\_physic/science-mango/phyx\_mini\_run/phyx\_data/test\_image/53.png

\paragraph{Formalization target.}
create a compiling Lean file with sorry bodies at `PhyXMiniProblems/problem\_phyx\_mini\_0053.lean`. The Lean declarations must preserve the physical quantities, dimensions or dimensional roles, figure labels, governing-law hypotheses, and final relation expressed by this problem.
Use Mathlib/Physlib names found through LeanExplore where available. If a physics API is missing, introduce faithful local abstractions rather than scalar placeholder aliases.

\begin{theorem}[Physics formalization target]
\label{thm:physics:phyx_mini_0053:target}
\lean{PhyXMiniProblems.ProblemPhyXMini0053.problem_phyx_mini_0053}
\uses{def:physics:phyx-mini-0053:phyxminiproblems-problemphyxmini0053-diverginglenssetup, def:physics:phyx-mini-0053:phyxminiproblems-problemphyxmini0053-matchesprimaryfigure, def:physics:phyx-mini-0053:phyxminiproblems-problemphyxmini0053-usescartesiansignconvention, def:physics:phyx-mini-0053:phyxminiproblems-problemphyxmini0053-hasphysicalsignconfiguration, def:physics:phyx-mini-0053:phyxminiproblems-problemphyxmini0053-obeyssignedthinlensequation, def:physics:phyx-mini-0053:phyxminiproblems-problemphyxmini0053-obeystransversemagnificationlaw, def:physics:phyx-mini-0053:phyxminiproblems-problemphyxmini0053-answerchoice, def:physics:phyx-mini-0053:phyxminiproblems-problemphyxmini0053-isnearesthundredthreadout}
The assigned autoformalize agent should translate this physics problem into Lean declarations in the covered file, with theorem and lemma proof bodies written as `by sorry`.
\end{theorem}
\begin{proof}
This is an autoformalization task, not a proof task. Produce statements that can later be proved by the physics prover without weakening the physical model.
\end{proof}

% --- Archon declaration topology begin ---
\section{Declaration topology}
\begin{definition}[Length Quantity]
  \label{def:physics:phyx-mini-0053:phyxminiproblems-problemphyxmini0053-lengthquantity}
  \lean{PhyXMiniProblems.ProblemPhyXMini0053.LengthQuantity}
  A signed physical length whose readout changes coherently with unit choice
\end{definition}
\begin{proof}
  This is the declared model definition of Length Quantity.
\end{proof}
\begin{definition}[centimeter Unit Choices]
  \label{def:physics:phyx-mini-0053:phyxminiproblems-problemphyxmini0053-centimeterunitchoices}
  \lean{PhyXMiniProblems.ProblemPhyXMini0053.centimeterUnitChoices}
  SI base units with centimeters selected as the length unit
\end{definition}
\begin{proof}
  This is the declared model definition of centimeter Unit Choices.
\end{proof}
\begin{definition}[length In Centimeters]
  \label{def:physics:phyx-mini-0053:phyxminiproblems-problemphyxmini0053-lengthincentimeters}
  \lean{PhyXMiniProblems.ProblemPhyXMini0053.lengthInCentimeters}
  \uses{def:physics:phyx-mini-0053:phyxminiproblems-problemphyxmini0053-lengthquantity, def:physics:phyx-mini-0053:phyxminiproblems-problemphyxmini0053-centimeterunitchoices}
  The signed scalar readout of a physical length in centimeters
\end{definition}
\begin{proof}
  This is the declared model definition of length In Centimeters.
\end{proof}
\begin{definition}[Thin Lens Kind]
  \label{def:physics:phyx-mini-0053:phyxminiproblems-problemphyxmini0053-thinlenskind}
  \lean{PhyXMiniProblems.ProblemPhyXMini0053.ThinLensKind}
  The two paraxial thin-lens types, distinguished by focal-length sign
\end{definition}
\begin{proof}
  This is the declared model definition of Thin Lens Kind.
\end{proof}
\begin{definition}[Geometrical Image Kind]
  \label{def:physics:phyx-mini-0053:phyxminiproblems-problemphyxmini0053-geometricalimagekind}
  \lean{PhyXMiniProblems.ProblemPhyXMini0053.GeometricalImageKind}
  Whether image rays meet physically or only after backward extension
\end{definition}
\begin{proof}
  This is the declared model definition of Geometrical Image Kind.
\end{proof}
\begin{definition}[Image Orientation]
  \label{def:physics:phyx-mini-0053:phyxminiproblems-problemphyxmini0053-imageorientation}
  \lean{PhyXMiniProblems.ProblemPhyXMini0053.ImageOrientation}
  The transverse orientation of an image relative to its object
\end{definition}
\begin{proof}
  This is the declared model definition of Image Orientation.
\end{proof}
\begin{definition}[Principal Axis Orientation]
  \label{def:physics:phyx-mini-0053:phyxminiproblems-problemphyxmini0053-principalaxisorientation}
  \lean{PhyXMiniProblems.ProblemPhyXMini0053.PrincipalAxisOrientation}
  The orientation of the principal optical axis in the source diagram
\end{definition}
\begin{proof}
  This is the declared model definition of Principal Axis Orientation.
\end{proof}
\begin{definition}[Figure Point]
  \label{def:physics:phyx-mini-0053:phyxminiproblems-problemphyxmini0053-figurepoint}
  \lean{PhyXMiniProblems.ProblemPhyXMini0053.FigurePoint}
  Labeled points whose signed axial positions occur in the ray diagram
\end{definition}
\begin{proof}
  This is the declared model definition of Figure Point.
\end{proof}
\begin{definition}[Figure Ray Label]
  \label{def:physics:phyx-mini-0053:phyxminiproblems-problemphyxmini0053-figureraylabel}
  \lean{PhyXMiniProblems.ProblemPhyXMini0053.FigureRayLabel}
  The two ray labels printed as a and b in the primary figure
\end{definition}
\begin{proof}
  This is the declared model definition of Figure Ray Label.
\end{proof}
\begin{definition}[Diverging Lens Setup]
  \label{def:physics:phyx-mini-0053:phyxminiproblems-problemphyxmini0053-diverginglenssetup}
  \lean{PhyXMiniProblems.ProblemPhyXMini0053.DivergingLensSetup}
  \uses{def:physics:phyx-mini-0053:phyxminiproblems-problemphyxmini0053-lengthquantity, def:physics:phyx-mini-0053:phyxminiproblems-problemphyxmini0053-thinlenskind, def:physics:phyx-mini-0053:phyxminiproblems-problemphyxmini0053-geometricalimagekind, def:physics:phyx-mini-0053:phyxminiproblems-problemphyxmini0053-imageorientation, def:physics:phyx-mini-0053:phyxminiproblems-problemphyxmini0053-principalaxisorientation, def:physics:phyx-mini-0053:phyxminiproblems-problemphyxmini0053-figurepoint, def:physics:phyx-mini-0053:phyxminiproblems-problemphyxmini0053-figureraylabel}
  The complete physical setup represented in the question and primary image. focalLengthMagnitude is positive, while signedFocalLength uses the Cartesian sign convention. objectDistance is the positive lens-to-object distance and signedImageDistance is negative for a virtual image on the object side. Heights are positive magnitudes because the pictured virtual image is upright
\end{definition}
\begin{proof}
  This is the declared model definition of Diverging Lens Setup.
\end{proof}
\begin{definition}[Matches Primary Figure]
  \label{def:physics:phyx-mini-0053:phyxminiproblems-problemphyxmini0053-matchesprimaryfigure}
  \lean{PhyXMiniProblems.ProblemPhyXMini0053.MatchesPrimaryFigure}
  \uses{def:physics:phyx-mini-0053:phyxminiproblems-problemphyxmini0053-lengthincentimeters, def:physics:phyx-mini-0053:phyxminiproblems-problemphyxmini0053-diverginglenssetup}
  Problem and figure readouts. The primary image places the object 100 cm left of the lens and the two focal points 50 cm from its center. It depicts an upright, reduced virtual image between the left focal point and the lens, and includes both rays labeled a and b. The circled 4, 5, and 6 in the image are callout markers, not numerical height or distance measurements, so no values are assigned to the two heights or to the image location here
\end{definition}
\begin{proof}
  This is the declared model definition of Matches Primary Figure.
\end{proof}
\begin{definition}[Uses Cartesian Sign Convention]
  \label{def:physics:phyx-mini-0053:phyxminiproblems-problemphyxmini0053-usescartesiansignconvention}
  \lean{PhyXMiniProblems.ProblemPhyXMini0053.UsesCartesianSignConvention}
  \uses{def:physics:phyx-mini-0053:phyxminiproblems-problemphyxmini0053-lengthincentimeters, def:physics:phyx-mini-0053:phyxminiproblems-problemphyxmini0053-diverginglenssetup}
  The Cartesian bookkeeping relating the signed distances used by the imaging laws to the positive focal-length magnitude and to the axial figure positions. This states a sign convention, not the requested image distance or magnification
\end{definition}
\begin{proof}
  This is the declared model definition of Uses Cartesian Sign Convention.
\end{proof}
\begin{definition}[Has Physical Sign Configuration]
  \label{def:physics:phyx-mini-0053:phyxminiproblems-problemphyxmini0053-hasphysicalsignconfiguration}
  \lean{PhyXMiniProblems.ProblemPhyXMini0053.HasPhysicalSignConfiguration}
  \uses{def:physics:phyx-mini-0053:phyxminiproblems-problemphyxmini0053-lengthincentimeters, def:physics:phyx-mini-0053:phyxminiproblems-problemphyxmini0053-diverginglenssetup}
  Qualitative sign conditions for the depicted diverging-lens configuration
\end{definition}
\begin{proof}
  This is the declared model definition of Has Physical Sign Configuration.
\end{proof}
\begin{definition}[Obeys Signed Thin Lens Equation]
  \label{def:physics:phyx-mini-0053:phyxminiproblems-problemphyxmini0053-obeyssignedthinlensequation}
  \lean{PhyXMiniProblems.ProblemPhyXMini0053.ObeysSignedThinLensEquation}
  \uses{def:physics:phyx-mini-0053:phyxminiproblems-problemphyxmini0053-diverginglenssetup}
  The paraxial signed thin-lens equation 1/f = 1/s + 1/s', expressed in the dimensionally homogeneous form f * (s + s') = s * s' in every unit system
\end{definition}
\begin{proof}
  This is the declared model definition of Obeys Signed Thin Lens Equation.
\end{proof}
\begin{definition}[Obeys Transverse Magnification Law]
  \label{def:physics:phyx-mini-0053:phyxminiproblems-problemphyxmini0053-obeystransversemagnificationlaw}
  \lean{PhyXMiniProblems.ProblemPhyXMini0053.ObeysTransverseMagnificationLaw}
  \uses{def:physics:phyx-mini-0053:phyxminiproblems-problemphyxmini0053-diverginglenssetup}
  The signed transverse-magnification laws m = -s'/s and h\_i = m h\_o. The distance-ratio equation determines the requested dimensionless quantity; the height equation records its physical interpretation in the ray diagram
\end{definition}
\begin{proof}
  This is the declared model definition of Obeys Transverse Magnification Law.
\end{proof}
\begin{definition}[Answer Choice]
  \label{def:physics:phyx-mini-0053:phyxminiproblems-problemphyxmini0053-answerchoice}
  \lean{PhyXMiniProblems.ProblemPhyXMini0053.AnswerChoice}
  Labels of the four displayed multiple-choice answers
\end{definition}
\begin{proof}
  This is the declared model definition of Answer Choice.
\end{proof}
\begin{definition}[magnification]
  \label{def:physics:phyx-mini-0053:phyxminiproblems-problemphyxmini0053-answerchoice-magnification}
  \lean{PhyXMiniProblems.ProblemPhyXMini0053.AnswerChoice.magnification}
  \uses{def:physics:phyx-mini-0053:phyxminiproblems-problemphyxmini0053-answerchoice}
  The dimensionless decimal magnification printed beside each answer choice
\end{definition}
\begin{proof}
  This is the declared model definition of magnification.
\end{proof}
\begin{definition}[Is Nearest Hundredth Readout]
  \label{def:physics:phyx-mini-0053:phyxminiproblems-problemphyxmini0053-isnearesthundredthreadout}
  \lean{PhyXMiniProblems.ProblemPhyXMini0053.IsNearestHundredthReadout}
  A displayed value is the nearest-hundredth readout of an exact magnification
\end{definition}
\begin{proof}
  This is the declared model definition of Is Nearest Hundredth Readout.
\end{proof}
\begin{lemma}[signed Image Distance In Centimeters eq neg hundred div three]
  \label{lem:physics:phyx-mini-0053:phyxminiproblems-problemphyxmini0053-signedimagedistanceincentimeters-eq-neg-hundred-div-three}
  \lean{PhyXMiniProblems.ProblemPhyXMini0053.signedImageDistanceInCentimeters_eq_neg_hundred_div_three}
  \uses{def:physics:phyx-mini-0053:phyxminiproblems-problemphyxmini0053-lengthincentimeters, def:physics:phyx-mini-0053:phyxminiproblems-problemphyxmini0053-diverginglenssetup, def:physics:phyx-mini-0053:phyxminiproblems-problemphyxmini0053-matchesprimaryfigure, def:physics:phyx-mini-0053:phyxminiproblems-problemphyxmini0053-usescartesiansignconvention, def:physics:phyx-mini-0053:phyxminiproblems-problemphyxmini0053-hasphysicalsignconfiguration, def:physics:phyx-mini-0053:phyxminiproblems-problemphyxmini0053-obeyssignedthinlensequation}
  The signed thin-lens equation with f = -50 cm and s = 100 cm determines the virtual-image distance s' = -100/3 cm
\end{lemma}
\begin{proof}
  The conclusion follows from the governing relations and figure readouts named in the statement of signed Image Distance In Centimeters eq neg hundred div three.
\end{proof}
\begin{lemma}[magnification eq one third]
  \label{lem:physics:phyx-mini-0053:phyxminiproblems-problemphyxmini0053-magnification-eq-one-third}
  \lean{PhyXMiniProblems.ProblemPhyXMini0053.magnification_eq_one_third}
  \uses{def:physics:phyx-mini-0053:phyxminiproblems-problemphyxmini0053-diverginglenssetup, def:physics:phyx-mini-0053:phyxminiproblems-problemphyxmini0053-matchesprimaryfigure, def:physics:phyx-mini-0053:phyxminiproblems-problemphyxmini0053-usescartesiansignconvention, def:physics:phyx-mini-0053:phyxminiproblems-problemphyxmini0053-hasphysicalsignconfiguration, def:physics:phyx-mini-0053:phyxminiproblems-problemphyxmini0053-obeyssignedthinlensequation, def:physics:phyx-mini-0053:phyxminiproblems-problemphyxmini0053-obeystransversemagnificationlaw}
  The transverse-magnification law and the derived virtual-image distance give the exact upright magnification m = 1/3
\end{lemma}
\begin{proof}
  The conclusion follows from the governing relations and figure readouts named in the statement of magnification eq one third.
\end{proof}
% --- Archon declaration topology end ---
% --- Archon physics formalization source end ---
